% !TEX root = 00_instructivo.tex
\chapter{Transductores eléctricos (cambios de resistencia), térmicos y de radiación}

\section*{Objetivos}
\begin{itemize}
    \item Medir humedad de sustrato (FC-28) y humedad ambiental (BME280) en un terrario.
    \item Evaluar la radiación UV con el sensor VEML6070.
    \item Comparar la precisión de un termistor (BMP280) frente a una termocupla tipo K.
\end{itemize}

\section*{Materiales y equipo}
\begin{itemize}
    \item Terrario con arena y sustrato.
    \item Sensor de humedad para sustrato FC-28.
    \item Sensor BME280 (humedad ambiental y temperatura barométrica).
    \item Sensor UV VEML6070.
    \item Sensor BMP280 (termistor).
    \item Termocupla tipo K con módulo MAX6675 (SPI).
    \item Arduino Uno.
    \item Mini protoboard de 9 canales.
    \item Cables Dupont (hembra-macho y macho-macho).
    \item Fuente USB y computadora con IDE Arduino.
\end{itemize}

\section*{Procedimiento}
\begin{enumerate}
    \item Conecte los sensores al Arduino según el esquemático.
    \item Cargue el código que lee todos los sensores periódicamente (cada 2 s).
    \item Coloque el FC-28 en el sustrato a 2 cm de profundidad.
    \item Registre durante 30 minutos: humedad del suelo, humedad ambiental, temperatura (BMP280 y termocupla), índice UV.
    \item Compare los valores de temperatura aplicando calor suave (lámpara incandescente) y ventilando.
    \item Elabore una tabla con promedios y desviaciones.
\end{enumerate}

\section*{Esquemático}
\begin{figure}[htbp]
\centering
\begin{circuitikz}[scale=0.8, transform shape]
\draw (0,0) node[above]{Arduino Uno};
% Pines de alimentación
\draw (2,0) node[ground]{GND};
\draw (2,0.5) node[anchor=south]{5V};
% Sensores I2C (BME280, VEML6070, BMP280) van a A4(SDA) y A5(SCL)
\draw (1,1) node[anchor=east]{A4 (SDA)} -- (3,1);
\draw (1,1.5) node[anchor=east]{A5 (SCL)} -- (3,1.5);
% FC-28: VCC a 5V, GND a GND, AO a A0
\draw (1,2) node[anchor=east]{A0} -- (3,2) node[anchor=west]{FC-28 (AO)};
% Termocupla MAX6675: SO a D12, CS a D10, SCK a D13
\draw (1,2.5) node[anchor=east]{D12 (SO)} -- (3,2.5);
\draw (1,3) node[anchor=east]{D10 (CS)} -- (3,3);
\draw (1,3.5) node[anchor=east]{D13 (SCK)} -- (3,3.5);
% Etiquetas
\node at (4,1) {BME280/VEML6070/BMP280 (I2C)};
\node at (4,2) {FC-28};
\node at (4,3) {MAX6675};
\end{circuitikz}
\caption{Conexiones en protoboard de 9 canales. Use el bus de 5V y GND para todos los sensores.}
\end{figure}

\textbf{Nota:} El BMP280 y el BME280 comparten el bus I2C (direcciones diferentes). El VEML6070 también es I2C.

\section*{Preguntas de análisis}
\begin{enumerate}
    \item ¿Qué diferencias observa entre las temperaturas del BMP280 y la termocupla? ¿A qué se deben?
    \item ¿Cómo influye la humedad del sustrato en la humedad ambiental dentro del terrario?
    \item Explique por qué el sensor UV requiere comunicación I2C y no una entrada analógica simple.
    \item Proponga una mejora para reducir el tiempo de respuesta de la termocupla.
\end{enumerate}

\section*{Bibliografía}
\begin{itemize}
    \item Hoja de datos FC-28. Rambal Electronics.
    \item Adafruit VEML6070 UV Sensor Breakout – Learning System.
    \item BME280 Datasheet. Bosch Sensortec.
    \item MAX6675 Datasheet. Maxim Integrated.
\end{itemize}
